% !TeX root = elegantnote-cn.tex
\section{化学动力学基础（一）}

\subsection{化学动力学的任务和目的 20260506}

\begin{enumerate}
    \item 热力学只能预言在给定的条件下，反应发生的可能性。至于如何把可能性变为现实性，以及过程中进行的速率如何，历程如何，热力学不能给出回答。
    \item 化学动力学的基本任务之一就是要了解反应的速率，了解各种因素（如分子结构、温度、压力、浓度、介质、催化剂等）对反应速率的影响。
    \item 化学动力学的另一个基本任务是研究反应历程，就是反应物究竞按什么途径、经过哪些步骤才转化为最终产物。
    \item 19世纪后半叶的宏观动力学阶段，主要成就是质量作用定律和 Arrhenius 公式的确立，并由此提出了活化能的概念。
    
    20世纪50年代以后的微观反应动力学阶段，主要是对反应速率从理论上进行了探讨，提出了碰撞理论和过渡态理论，并借助于量子力学计算了反应系统的势能面，指出过渡态（或活化络合物）乃是势能面上的马鞍点。

    近百年来化学动力学进展的速度很快，这一方面应归功于相邻学科基础理论和技术上的进展，另一方面归功于实验方法、检测手段的日新月异。 
\end{enumerate}

\subsection{化学反应速率的表示法 20260506}

\begin{enumerate}
    \item 反应速率表示浓度随时间的变化率。由于产物和反应物的化学计量数不尽一致，固用其浓度变化率来表示反应速率时，数值未必一致。
    \item 采用反应进度 $\xi$ 随时间的变化率来表示反应速率。设反应为
    $$\begin{array}{c c c c}
        & \ch{$\alpha$ R} & \ch{->} & \ch{$\beta$ P} \\
        t = 0 & n_{\ch{R}}(0) & & n_{\ch{P}}(0) \\
        t = t & n_{\ch{R}}(t) & & n_{\ch{P}}(t) \\
    \end{array}$$
    则反应进度为
    \begin{equation}
        \xi = \frac{n_{\ch{R}}(t)-n_{\ch{R}}(0)}{-\alpha} = \frac{n_{\ch{P}}(t)-n_{\ch{P}}(0)}{\beta}
    \end{equation}
    对 $t$ 微分，得到反应的转化速率
    \begin{equation}
        \frac{\mathrm{d}\xi}{\mathrm{d}t} = \dot{\xi} 
        = -\frac1{\alpha}\frac{\mathrm{d}n_{\ch{R}}(t)}{\mathrm{d}t}
        = \frac1{\beta} \frac{\mathrm{d}n_{\ch{P}}(t)}{\mathrm{d}t}
    \end{equation}

    \item 化学反应速率 $r$ 定义为
    \begin{equation}
        r \xlongequal{\text{def}} \frac1{V} \frac{\mathrm{d}\xi}{\mathrm{d}t} 
        = -\frac1{V \alpha} \frac{\mathrm{d}n_{\ch{R}}(t)}{\mathrm{d}t} = \frac1{V \beta} \frac{\mathrm{d}n_{\ch{P}}(t)}{\mathrm{d}t}
        \label{ReactionRate}
    \end{equation}
    式中 $V$ 为反应系统的体积。若反应过程中体积恒定，则式 \eqref{ReactionRate} 可写为
    \begin{equation*}
        \begin{aligned}
            r &= -\frac1{\alpha} \frac{\mathrm{d}[n_{\ch{R}}(t)/V]}{\mathrm{d}t} 
            = -\frac1{\alpha} \frac{\mathrm{d}c_{\ch{R}}}{\mathrm{d}t}
            = -\frac1{\alpha} \frac{\mathrm{d}[\ch{R}]}{\mathrm{d}t}\\
            &= \frac1{\beta} \frac{\mathrm{d}[n_{\ch{P}}(t)/V]}{\mathrm{d}t}
            = \frac1{\beta} \frac{\mathrm{d}c_{\ch{P}}}{\mathrm{d}t}
            = \frac1{\beta} \frac{\mathrm{d}[\ch{P}]}{\mathrm{d}t} \\
        \end{aligned}
    \end{equation*}

    \item 对于任意反应，有
    \begin{equation}
        r = \frac1{\nu_{\ch{B}}} \frac{\mathrm{d}[\ch{B}]}{\mathrm{d}t}
    \end{equation}
    式中 $r$ 的量纲为 $[\text{浓度}] \cdot [\text{时间}]^{-1}$。

    \item 气相反应，可用参加反应各物种的分压代替浓度。
    \begin{equation}
        r' = \frac1{\nu_{\ch{B}}} \frac{\mathrm{d}p_{\ch{B}}}{\mathrm{d}t}
    \end{equation}
    式中 $r'$ 的量纲为 $[\text{压力}] \cdot [\text{时间}]^{-1}$。对于理想气体，$p_{\ch{B}} = c_{\ch{B}} RT$，故 $r' = r \cdot RT$。

    \item 对于多相催化反应，反应速率可定义为
    \begin{equation}
        r \xlongequal{\text{def}} \frac1{Q} \frac{\mathrm{d}\xi}{\mathrm{d}t}
    \end{equation}
    \begin{enumerate}
        \item $Q$ 用质量 $m$ 表示：$r_m = \dfrac1{m} \dfrac{\mathrm{d}\xi}{\mathrm{d}t}$，称为催化剂的比活性，单位为 \si{\mole\per\kilo\gram\per\second}。 \\
        \item $Q$ 用堆体积 $V$ 表示：$r_V = \dfrac1{V} \dfrac{\mathrm{d}\xi}{\mathrm{d}t}$，单位为 \si{\mole\per\meter\cubed\per\second}。 \\
        \item $Q$ 用表面积 $A$ 表示：$r_A = \dfrac1{A} \dfrac{\mathrm{d}\xi}{\mathrm{d}t}$，称为表面反应速率，单位为 \si{\mole\per\meter\squared\per\second}。 \\
    \end{enumerate}

    \item 绘制动力学曲线：反应中各物质浓度随时间的变化曲线。在时间 $t$ 时作曲线的切线，可求反应速率。
    \item 测定物质不同时刻浓度的方法：
    \begin{enumerate}
        \item 化学方法：不同时刻取出一定量反应物，设法（骤冷、稀释、阻化、去催化）使反应立即停止，然后化学分析。
        \item 物理方法：用各种方法测定与浓度有关的物理性质（压力、体积、旋光度、折射率、吸收光谱、电导、电动势、介电常数、黏度、热导率或进行比色），或用现代谱仪检测与浓度有定量关系的物理量的变化，求得浓度变化。可获得原位反应数据。
        \item 反应时间很短的反应采用特殊装置测量，如快速流动法。
    \end{enumerate}
\end{enumerate}
%----------%
\subsection{化学反应的速率方程 20260506 20260511}

\begin{enumerate}
    \item 表示反应速率与浓度等参数之间的关系的方程或表示浓度等参数与时间关系的方程称为化学反应的速率方程，也称为动力学方程。
    \item 速率方程可表示为微分式或积分式，必须由实验来测定。
\end{enumerate}

\subsubsection{基元反应和非基元反应}

\begin{enumerate}   
    \item 分子经一次碰撞后，在一次化学行为中就能完成的反应称为\textbf{基元反应}。
    \item \textbf{反应机理}
    \item 质量作用定律：基元反应的祖率与反应物浓度（含有相应的整数）的乘积成正比，其中各浓度的级数
\end{enumerate}

\subsubsection{反应的级数、反应分子数和反应的速率常数}

\begin{enumerate}
    \item 各反应浓度项上的指数称为称为该反应物的级数
    \item 实验测定
    \item 在基元反应中，实际参加反应的分子数称为反应分子数。对于基元反应或简单反应，通常其反应级数和反应的分子数是相同的。
    \item 但也有些基元反应表现出的反应级数与反应分子数不一致。
    \item 通常基元反应的级数为不超过3。
    \item 在式子（）中比例系数 $k$ 与浓度无关，称为速率常数或速率系数。
    \item 不同反应的有不同的速率常数，速率常数与反应温度、反应介质、催化剂有关。
    \item $k$ 的数值反映了速率的快慢。
\end{enumerate}

\subsection{具有简单级数的反应 20260511 20260513}

\subsubsection{一级反应}

\begin{enumerate}
    \item 反应速率知鱼反应物浓度的一次方成正比。
    \item 放射性元素、同位素的衰变，分子重排反应，蔗糖水解反应。
    \item 设某一级反应
    \[
    \begin{array}{c c c c}
        & \ch{A} & \ch{->[ $k_1$ ]} & \ch{E} \\
        t = 0 & c_{\mathrm{A}}^0 & & c_{\mathrm{P}}^0 \\
        t = t & c_{\mathrm{A}} = a - x & & c_{\mathrm{P}} = x \\
    \end{array}
    \]

    \begin{equation}
        \begin{aligned}
            r = -\frac{\mathrm{d}c_{\mathrm{A}}}{\mathrm{d}t}
        \end{aligned}
    \end{equation}
    \item 一级反应的性质
    
    $\ln(a-x)$ 与 $t$ 成正比；
    
    $k_1 = \frac1t \ln{\frac{a}{a-x}}$，速率常数的单位是时间的负一次方；
    
    $\frac{a-x}{a} = \exp{-k_1t}$，时间无限延长，反应物浓度趋于零；
    
    令 $y = \frac{x}{a}$，（反应物的转化率），带入速率方程，$\ln\frac1{1-y} = k_1 t$，当 $y = \frac12$，对应的时间为半衰期 $t = \frac{\ln2}{k_1}$。
    \item 引申特点
\end{enumerate}

\subsubsection{二级反应}

\begin{enumerate}

    \item 反应速率和物质浓度的二次方成正比。
    \item 第一种
    \begin{enumerate}
        \item 
        \begin{equation}
            \frac1{a-x} - \frac1a = k_2 t 
        \end{equation}
        \item 特点
    
        1/(a-x)与t成正比

        $k_2 = \frac1t \frac{x}{a(a-x)}$，速率常数的单位是时间和浓度的负一次方；

        $y = \frac{x}{a}$，$\frac{y}{1-y} = k_2 t a$，$y = \frac12$ 时，半衰期……
        \item 引申特点
        \item $a \neq b$ 时，……，没有统一半衰期表示式
    \end{enumerate}
    \item 第二种
    \begin{enumerate}
        \item dsd
        \begin{equation}
            \frac{\mathrm{d}x}{\mathrm{d}t} = k_2 (a - 2x)^2
        \end{equation}
        \begin{equation}
            \frac{x}{a(a-2x)} = k_2 t
        \end{equation}
        \item 在二级反应中用浓度表示的速率常数和用压力表示的速率常数在树枝上不相等，而在一级反应中二者是相等的。
        \begin{equation}
            r = -\frac12 \frac{\mathrm{d}[A]}{\mathrm{d}t} = k_2 [A]^2
        \end{equation}
        \begin{equation}
            \frac{k_2}{RT} = k_p
        \end{equation}
    \end{enumerate}
\end{enumerate}



\subsubsection{三级反应}

\begin{enumerate}
    \item 反应速率与物质浓度的三次方成正比。
    \item 可分为以下三种反应：
    \begin{comment}
        A + B + C
        \arrow{->}
        P
    
    \end{comment}
       
    \schemestart
        2A + B
        \arrow{->}
        P
    \schemestop \\
    \schemestart
        3A
        \arrow{->}
        P
    \schemestop
    \item displ
    \begin{equation}
        \frac12 \left[\frac1{(a-x)^2} - \frac1{a^2}\right] = k_3 t
    \end{equation}
    \begin{equation}
        t_{1/2} = \frac3{2k_3 a^2}
    \end{equation}
    \item 三级反应的特点
\end{enumerate}

\subsubsection{零级反应和准级反应}

\begin{enumerate}
    \item 零级反应
    \begin{enumerate}
        \item 反应速率与物质的浓度无关。
        \item E
        \begin{equation}
            r = - \frac{\mathrm{d}c_{\mathrm{A}}}{\mathrm{d}t} = k_0 \quad \text{或} \quad  r = - \frac{\mathrm{d}x}{\mathrm{d}t} = k_0
        \end{equation}
        移项积分，得
        \begin{equation}
            x = k_0 t
        \end{equation}
        当 $\displaystyle x = \frac{a}{2}$ 时，$\displaystyle t_{1/2} = \frac{a}{2k_0}$。
        \item 零级反应的特点
    \end{enumerate}
    \item 准级反应
    \begin{enumerate}
        \item 某一物质的浓度远大于其他反应物的浓度，或是出现在速率方程中催化剂浓度项，再反应过程中可以认为没有变化，可并入速率常数项，反应级数降低。
    \end{enumerate}
    \item n 级反应
    \begin{enumerate}
        \item 仅有一种反应物，$n \neq 1$。 
        \item nA -> p \\
        t = 0 a 0 \\
        t = t a-x x
        \begin{equation}
            r = \frac{\mathrm{d}x}{\mathrm{d}t} = k(a-x)^n
        \end{equation}
        \begin{equation}
            \int_{0}^{x} \frac{\mathrm{d}x}{(a-x)^n} = \int_{0}^{t} k\mathrm{d}t
        \end{equation}
        \item 半衰期的一般式
        \begin{equation}
            t_{\mathrm{1/2}} = 
        \end{equation}
        \item n 级反应的特点
    \end{enumerate}
\end{enumerate}

\subsubsection{反应技术的测定法}

\begin{enumerate}
    \item 积分法（尝试法）
    \begin{enumerate}
        \item 假设 $\alpha$，$\beta$ 的值和反应级数，带入公式。若得到 $k$ 为一常数，则假设正确。要求掌握。
        \item 作图法
    \end{enumerate}
    \item 微分法
    \begin{enumerate}
        \item 用 n 级反应的速率方程求出反应级数。
        \begin{eqnarray}
            r = -\frac{\mathrm{d}c_{\mathrm{A}}}{\mathrm{d}t} =
        \end{eqnarray}
        \item 需要作3次图：$c_\mathrm{A} \sim t$，不同时刻求 $-\frac{\mathrm{d}c_{\mathrm{A}}}{\mathrm{d}t}$，$\ln\left(-\frac{\mathrm{d}c_{\mathrm{A}}}{\mathrm{d}t}\right) \sim \ln{c_{\mathrm{A}}}$
    \end{enumerate}
    \item 半衰期法
    \begin{enumerate}
        \item 反应物起始浓度相同，则
        \begin{equation}
            t_{\mathrm{1/2}} = A \frac1{a^{n-1}}
        \end{equation}
        \item 亦可使用作图法
    \end{enumerate}
    \item 改变物质数量比例的方法，类似于准级反应。
\end{enumerate}

\subsection{几种典型的复杂反应 20260513 20250518}

\subsubsection{对峙反应}

\begin{enumerate}
    \item 在正、反两个方向上都能进行的反应称为对峙反应，亦称为可逆反应。
    \item 正、逆反应可以为相同级数，也可以为具有不同级数的反应；可以是基元反应，也可以是非基元反应。
    \item $1-1$ 级对峙反应
    \schemestart
    A \arrow{<=>[$k_1$][$k_{-1}$]} B
    \schemestop
    \item $2-2$ 级对峙反应
    \schemestart
    A $+$ B \arrow{<=>[$k_1$][$k_{-1}$]} C $+$ D
    \schemestop
    \item 对峙反应的特点
\end{enumerate}

\subsubsection{平行反应}

\begin{enumerate}
    \item 两个反应都是一级反应
    \begin{equation}
        \ln\frac{a}{a-x} = (k_1 + k_2) t
    \end{equation}
    \item 两个反应都是二级反应
    \begin{equation}
        \frac1{a-x} - \frac1a = (k_1 + k_2) t
    \end{equation}
    \item 特点
    \begin{enumerate}
        \item 反应总速率等于各平行反应速率之和
        \item 速率常数为各平行反应速率常数之和
        \item 起始浓度为零，反应级数相等，任意时间各产物浓度等于速率常数之比。
        \item 用合适催化剂可改变某一反应的速率，提高主反应产物的含量。
        \item 改变温度也可改变产物的相对含量。活化能高的反应，温度系数随温度的变化率也大。
        \begin{equation*}
            \frac{\mathrm{d}\ln{k}}{\mathrm{d}T} = \frac{E_\mathrm{a}}{RT^2}
        \end{equation*}
        \begin{equation*}
            \frac{\mathrm{d}\ln{k_1 / k_2}}{\mathrm{d}T} = \frac{E_{\mathrm{a}1}-E_{\mathrm{a}2}}{RT^2}
        \end{equation*}
    \end{enumerate}   
\end{enumerate}

\subsubsection{连续反应}

\begin{enumerate}
    \item 经过连续几步才完成，前一步的产物就是下一步的反应物
    \item 两个连续一级反应
    \begin{equation*}
        \begin{aligned}
            x + y + z = a \\
            -\frac{\mathrm{d}x}{\mathrm{d}t} = k_1 x \quad \int_{a}^{x} \frac{\mathrm{d}x}{x} = \int_{0}^{t} k_1 \mathrm{d}t \quad x = a \mathrm{e}^{-k_1 t} \\
            \frac{\mathrm{d}y}{\mathrm{d}t} = k_1 x - k_2 y \quad y = \frac{k_1 a}{k_2 - k_1} (\mathrm{e}^{-k_1 t} - \mathrm{e}^{-k_2 t}) \\
            \frac{\mathrm{d}z}{\mathrm{d}t} = k_2 y \\
        \end{aligned}
    \end{equation*}
    \item 讨论
    \begin{enumerate}
        \item $k_1 \gg k_2$
        \begin{equation*}
            z = a \left(1 + \frac{k_2}{k_1} \mathrm{e}^{-k_1 t} - \frac{k_1}{k_1} \mathrm{e}^{-k_2 t}\right) = a \left(1 - \mathrm{e}^{-k_2 t}\right)
        \end{equation*}
        \item $k_1 \ll k_2$
        \begin{equation*}
            z = a \left(1 - \mathrm{e}^{-k_1 t}\right)
        \end{equation*}
    \end{enumerate}
    \item 连续反应的 $c - t$ 关系曲线：B 的浓度先增加后减少。 % pic
    B 浓度的最大值 $\frac{\mathrm{d}y}{\mathrm{d}t} = 0$
    \begin{equation*}
        t_{\mathrm{m}} = \frac{\ln{k_2} - \ln{k_1}}{k_2 - k_1}
    \end{equation*}
    \begin{equation}
        y_{\mathrm{m}} = a \left(\frac{k_1}{k_2}\right)^{\frac{k_2}{k_2 - k_1}}
    \end{equation}
\end{enumerate}

\subsection{基元反应的微观可逆性原理}

\subsection{温度对反应速率的影响 20260518}

\subsubsection{速率常数与温度的关系 —— Arrhenius 经验式}

\begin{enumerate}
    \item van't Hoff 近似规则：温度每升高 \qty{10}{\kelvin}，反应速率增加 $2 \sim 4$ 倍。
    \item Arrhenius 经验式提出了活化能的概念，并揭示了反应的速率常数与温度的依赖关系
    \begin{equation}
        k = A \mathrm{e}^{-\frac{E_\mathrm{a}}{RT}}
    \end{equation}
    $k$ 是温度为 $T$ 时反应的速率常数；$R$ 是摩尔气体常数，$A$ 是指前因子；$E_{\mathrm{a}}$ 是表观活化能，通常简称为活化能。
    \item 取对数、微分
    \item 热力学和动力学对温度对速率常数影响的看法
    
    吸热反应
    \begin{equation*}
        \begin{aligned}
            \frac{\mathrm{d}\ln{K^{\circ}}}{\mathrm{d}T} = \frac{\Delta_{\mathrm{r}} H_{\mathrm{m}}^{\circ}}{RT^2}
        \end{aligned}
    \end{equation*}
    放热反应出现矛盾

    速率因素是主要因素，因此可适当牺牲一些平衡转化率以提高反应速率。
\end{enumerate}

\subsubsection{反应速率与温度关系的几种类型}

\begin{enumerate}
    \item 
\end{enumerate}

\subsection{关于活化能}

\subsection{链反应 20260518}

\begin{enumerate}
    \item 链的开始/链的引发、链的传递/链的增长、链的终止
    \item 直链反应
    \begin{enumerate}
        \item 已知总包反应 $\ch{H2(g)} + \ch{Cl2(g)} \ch{->} \ch{2 HCl(g)}$
        \item 推测反应机理
        \item 速率方程
        \begin{equation*}
            r = \frac12 \frac{\mathrm{d}[\ch{HCl}]}{\mathrm{d}t} = k[\ch{Cl2}]^{\frac12}[\ch{H2}]
        \end{equation*}
    \end{enumerate}
\end{enumerate}

\subsubsection{直链反应（\ch{Cl2} 反应的历程) —— 稳态近似法}
令浓度随时间变化率，写出关于产物有关式子的速率方程，代入上式可抵消

\subsubsection{支链反应}